% Part: history
% Chapter: biographies 
% Section: emmy-noether

\documentclass[../../../include/open-logic-section]{subfiles}

\begin{document}

\olfileid{his}{bio}{noe}

\olsection{اېمي نوېتر}

\olphoto{noether-emmy}{اېمي نوېتر}

اېمي نوېتر (\textsc{نېړ}-ټر) د 1882 د مارچ پر
23مه د جرمني په اېرلانګن کښې د لوړې منځنۍ
طبقې په يوه علمي کورنۍ کښې زېږېدلې وه۔
نوېتر، چې «د نوي الجبر مور» بلل کېږي، د
ښځو د زده کړو له لويو خنډونو سره سره په
رياضياتو او فزيک دواړو کښې بنسټيزې نوې
مرستې وکړې۔ په هغه وخت کښې په جرمني کښې
ځوانو نجونو ته د هنرونو زده کړه ټاکل کېده
او د پوهنتون د چمتووالي ښوونځيو ته د
تللو اجازه ئې نۀ لرله۔ خو د ګوټينګن او
اېرلانګن په پوهنتونونو کښې د ټولګيو له
اورېدو وروسته، چې پلار ئې په اېرلانګن
کښې د رياضياتو پروفېسر و، نوېتر ته بالاخره
په 1904 کښې هلته د زده کوونکې په توګه
د نوم‌ليکنې شونتيا برابره شوه، کله چې د پوهنتون
تګلاره د ښځينه زده کوونکو د منلو دپاره
بدله شوه۔ هغې په 1907 کښې په رياضياتو
کښې دوکتورا واخيسته۔

له خپل اهليت سره سره نوېتر په خپل کاري
ژوند کښې له ډېر مخالفت سره مخ شوه۔ له
1908--1915 پورې ئې په اېرلانګن کښې بې
معاشه درس ورکاوۀ۔ په دې وخت کښې ئې د
ډېوېډ هېلبرټ پام ځان ته راواړوۀ، چې د
هغه وخت د نړۍ له مخکښو رياضيپوهانو څخه
و او هغه ئې ګوټينګن ته راوبلله۔ خو ښځو
ته د پروفېسر د مقام ترلاسه کول منع وو،
او هغې يوازې د هېلبرټ تر نامه لاندې
درس ورکولے شو، بيا هم بې معاشه۔ په دې
وخت کښې ئې هغه څه ثابت کړل چې اوس د
نوېتر قضيه نومېږي او نن هم په نظري
فزيک کښې کارېږي۔ نوېتر ته بالاخره په
1919 کښې د درس ورکولو حق ورکړے شو۔
ويل کېږي چې د پوهنتون د همکارانو د
جاري مخالفت په ځواب کښې هېلبرټ وويل:
«ښاغلو، د پوهنځي سنا حمام نۀ دے۔»

د 1920s په وروستيو کښې هغې په تجريدي الجبر
کار ته پام واړاوۀ، او مرستو ئې دا ډګر
واړوۀ۔ په خپلو ثبوتونو کښې ئې ډېر ځله
د لوړېدونکې لړۍ په نوم شرط کاراوۀ، چې
وايي د ځينو سټونو هېڅ نامتناهي سخته
زياتېدونکې لړۍ نشته۔ د بېلګې په توګه،
ځينې جبري جوړښتونه چې اوس د نوېتر
حلقې نومېږي، دا خاصيت لري چې د ايډيالونو
هېڅ نامتناهي لړۍ $I_1
\subsetneq I_2 \subsetneq \dots$ نشته۔
دا شرط هر جزوي ترتيب ته عامېدلے شي
(په الجبر کښې دا د فرعي سټ د اړيکې
له مخې د مرتب شوو ايډيالونو ځانګړے
حالت دے)، او موږ ئې دوه‌ګون، د
ښکته کېدونکې لړۍ شرط هم څېړلے شو،
چې په جزوي ترتيب کښې هره سخته
\emph{کمېدونکې} لړۍ بالاخره ختمېږي۔
که يو جزوي ترتيب د ښکته کېدونکې لړۍ
شرط پوره کوي، نو پر همدغه ترتيب
استقرا په هغه شان کارولے شو چې
د~$<$ ترتيب پر~$\Nat$ استقرا کاروو۔
داسې ترتيبونه \emph{ښه‌بنسټه} يا
\emph{د نوېتر ترتيبونه} نومېږي، او
اړوند ثبوتي اصل \emph{د نوېتر استقرا}
نومېږي۔

نوېتر يهودۍ وه، او کله چې نازيان په 1933
کښې واک ته ورسېدل، هغه له خپلې دندې
ګوښه شوه۔ له نېکه مرغه نوېتر په
پنسلوانيا کښې په برين ماور کښې د
يوې موقتي دندې دپاره متحده ايالاتو
ته کډواله کېدلے شوه۔ هلته د اوسېدو
پر وخت ئې په پرنسټن کښې هم درسونه
ورکول، که څه هم هغې وليدل چې پوهنتون
د ښځو هرکلے نۀ کاوۀ
\citep[81]{Dick1981}۔ په 1935 کښې
د نوېتر د رحم د تومور د لرې کولو
عمليات وشول۔ هغه د همدغو عملياتو
په پايله کښې د عفونت له امله وفات
شوه او په برين ماور کښې خاورو ته
وسپارل شوه۔

\begin{reading}
د نوېتر د ژوندليک دپاره \citet{Dick1981}
وګورئ۔ د نظري فزيک پېريميټر انسټيټيوټ
د نوېتر د ژوند او اغېز درسونه پر
انټرنېټ وړاندې کوي \citep{Perimeter2015}۔
که له لوستلو ستړي شوي يئ، \emph{هغه څه
چې د تاريخ په ټولګي کښې مو هېر کړي}
د نوېتر د ژوند او اغېز يو پوډکاسټ لري
\citep{Frey2015}۔ د نوېتر راټولې شوې
ليکنې په اصلي جرمني ژبه شته
\citep{Noether1983}۔
\end{reading}

\end{document}
